% =========================================================
%  INSERT 3 of 3:  New appendix app:udg_stress_test
% =========================================================
%  Proposed insertion location:
%  Between two existing appendices in the galactic-sector cluster:
%  AFTER  \section{ECT derivation status for the galactic interpolation law}
%         (line 51838, label: app:rar_derivation_status)
%         ... ends around line ~51902
%  BEFORE \section{Numerical algorithm for late-time ECT cosmology...}
%         (line 51904, label: app:late_cosmo_algorithm)
%
%  Rationale:
%    - logical flow is: derivation status -> stress test of that closure
%      -> cosmological algorithm (new topic)
%    - thematic cluster is preserved
%    - no existing hardcoded "Appendix X" references exist in the
%      preprint (verified via grep), so all \ref{app:...} commands
%      continue to resolve correctly after the insertion --- only
%      the letter labels shift automatically
%    - appendix letter order is thereby preserved in meaning
%      (galactic-sector-cluster-then-cosmology), not disturbed
%
%  Word count: ~2000 words; scientific content derivation + diagnostic
%  only.  Contains NO claim that ECT explains DF4, NO claim of derived
%  rescue mechanism, NO claim that bar-phi -> 0 is a derived limit.
% =========================================================

% ============================================================
\section{Diagnostic stress test of the current Level-B galactic
closure: dark-matter-deficient ultra-diffuse galaxies}
\label{app:udg_stress_test}

This appendix is a \emph{diagnostic appendix}: it applies the current
practical galactic closure
\eqref{eq:ect_rotation_closure_main} as an \emph{inverse-problem
probe} to a class of observed objects --- the dark-matter-deficient
ultra-diffuse galaxies (UDGs) --- and documents the resulting
stress test.  It does not propose a rescue mechanism and does not
claim that the current closure explains these objects.  Its purpose
is to record, in one place and with one consistent methodology, the
empirical tension between the current Level-B closure and the
low-dispersion UDG class, the accompanying nuisance budget, the
matched-pair observation, and the theory-level directions that a
future completion of ECT's galactic sector would need to address.

\textit{Status: diagnostic.  The scientific conclusions below are of
the form ``the current Level-B closure is strained on this object'',
not ``ECT predicts this observation''.  No rescue mechanism is
derived here.}

\subsection{Purpose and scope}
\label{app:udg:purpose}

The preprint's main galactic sector (\S\ref{sec:galactic_dynamics})
and the derivation-status appendix
(App.~\ref{app:rar_derivation_status}) establish that the current
practical closure \eqref{eq:ect_rotation_closure_main} has two
asymptotic branches --- Newtonian at large $g_N$ and deep-MOND at
small $g_N$ --- and that the minimal
ansatz~\eqref{eq:gdagger_env_main} for the environment modulator
$\Xi$ is an open element of the Level-B construction (OP3).  For
SPARC-like normal galaxies this closure reproduces the observed
radial-acceleration phenomenology at the level documented in
App.~\ref{app:galactic_phi_closure} and
App.~\ref{app:sparc_computation}.

Three objects of the DM-deficient UDG class --- NGC~1052--DF4
\cite{Danieli2022DF4,vanDokkum2018}, FCC~224
\cite{Buzzo2025FCC224,Tang2025FCC224}, and the observationally
ambiguous NGC~1052--DF2
\cite{Trujillo2019DF2,Shen2021DF2TRGB,vanDokkum2018} ---
have line-of-sight velocity dispersions consistent with
baryon-only Newtonian mass within their probed radii.  A
matched-mass DM-rich control (NGC~5846-UDG1
\cite{Forbes2021NGC5846UDG1}) lies on the standard RAR.  This
appendix quantifies the resulting inverse-problem diagnostic.

\subsection{Inverse-problem closed form: $\Xi_{\rm req}$ from
$\sigma_{\rm obs}$}
\label{app:udg:inverse_closed_form}

At a fixed radius $R$, the practical closure
\eqref{eq:ect_rotation_closure_main} reads
\begin{equation}
g(R)=\tfrac{1}{2}\!\left[g_N(R)+\sqrt{g_N^2(R)+4g_N(R)\,g^\dagger_{\rm eff}}\right],
\qquad
g^\dagger_{\rm eff}=g^\dagger_0\,\Xi,
\label{eq:app_udg_closure}
\end{equation}
with $g^\dagger_0=cH_0/(2\pi)\simeq 1.04\times 10^{-10}\,\text{m s}^{-2}$.
Define the local dynamical/baryonic ratio
\begin{equation}
r(R)\equiv\frac{M_{\rm dyn}(<R)}{M_{\rm bar}(<R)}
=\frac{g(R)}{g_N(R)}.
\end{equation}
Solving \eqref{eq:app_udg_closure} for $r$ gives
\begin{equation}
r=\tfrac{1}{2}\!\left[1+\sqrt{1+4y}\right],
\qquad y\equiv\frac{g^\dagger_{\rm eff}}{g_N(R)},
\end{equation}
whence the clean inversion
\begin{equation}
\boxed{\;
y=r(r-1),
\qquad
\Xi_{\rm req}(R)
=\frac{r(r-1)\,g_N(R)}{g^\dagger_0}.
\;}
\label{eq:app_udg_xi_req}
\end{equation}
This is a closed-form inverse of
\eqref{eq:app_udg_closure}.  For a pressure-supported system the
right-hand side is evaluated at the Wolf half-mass radius
$R_{1/2}\simeq\tfrac{4}{3}R_e$, where the isotropic Wolf estimator
\cite{Wolf2010} gives
\begin{equation}
M_{\rm dyn,1/2}=4\sigma_{\rm los}^2\,R_{1/2}/G,
\qquad
M_{\rm bar,1/2}\simeq\tfrac{1}{2}M_\star,
\label{eq:app_udg_wolf}
\end{equation}
so that $r_{\rm obs}=\sigma_{\rm los,obs}^2/\sigma_N^2$ with
$\sigma_N^2\equiv (3/32)GM_\star/R_e$ the Newtonian Wolf proxy for
isotropic Plummer.  Formula \eqref{eq:app_udg_xi_req} is \emph{not}
a prediction of $\Xi$; it is the value of $\Xi$ that the current
closure would need at this object to reproduce the observed
kinematics, given proxy assumptions \eqref{eq:app_udg_wolf}.

The remaining closure-level phenomenology is limited to: (i)~the
Wolf coefficient $4$, valid for isotropic Plummer with $\sim 30\%$
uncertainty in the proxy relation; (ii)~the half-mass fraction
$f_e=0.5$ for Plummer, varying by $\sim 20\%$ across plausible
S\'ersic indices.  Reducing these requires a full Jeans forward
model~(\S\ref{app:udg:jeans_upgrade}).

\subsection{Jeans upgrade path (phenomenology reduction, not elimination)}
\label{app:udg:jeans_upgrade}

A stricter treatment replaces \eqref{eq:app_udg_wolf} by the
isotropic spherical Jeans equation
\begin{equation}
\sigma_r^2(r)=\frac{1}{\nu(r)J_\beta(r)}
\int_r^\infty\!\nu(s)\,g(s)\,J_\beta(s)\,ds,
\qquad
J_\beta(r)=\exp\!\left(\int^r\!\frac{2\beta(u)}{u}du\right),
\end{equation}
followed by line-of-sight projection and aperture averaging,
\begin{equation}
\sigma_{\rm los}^2(R)=\frac{2}{I(R)}
\int_R^\infty\!\!\left(1-\beta\frac{R^2}{r^2}\right)
\frac{\nu(r)\sigma_r^2(r)\,r\,dr}{\sqrt{r^2-R^2}},
\qquad
\langle\sigma_{\rm los}^2\rangle_{\rm ap}=
\frac{\int_0^{R_{\rm ap}}\!W I\,\sigma_{\rm los}^2\,R\,dR}
{\int_0^{R_{\rm ap}}\!W I\,R\,dR},
\end{equation}
with $\nu(r)$ the deprojected tracer density from the observed
photometry, $\beta(r)$ the anisotropy, and $g(r)$ from
\eqref{eq:app_udg_closure} at trial $\Xi$.  For DF4 under isotropic
Plummer, $M/L=2$, and $D=20$~Mpc, a direct Jeans integration
returns $\sigma_N^{\rm Jeans}\simeq 8.0$~km\,s$^{-1}$ against the
Wolf-proxy value $\sigma_N^{\rm proxy}\simeq 6.2$~km\,s$^{-1}$
(ratio $1.30$).  The qualitative conclusion --- that DF4 requires
$\Xi_{\rm req}\ll 1$ under the current closure --- is unchanged.
The Jeans upgrade reduces but does not eliminate the residual
phenomenology, which still contains the observational inputs
$M/L$, $\beta(r)$, tracer family, aperture, and interloper
fraction.

\subsection{Numerical stress-test diagnostic}
\label{app:udg:numerical}

Applying \eqref{eq:app_udg_xi_req} at $R_{1/2}=\tfrac{4}{3}R_e$ to the
six reference objects gives Table~\ref{tab:udg_stress_test}.  The
quoted $\Xi_{\rm req}$ values are central diagnostic proxy-level
estimates; the full 95\% CL bands are summarised separately
in~\S\ref{app:udg:nuisance}.

\begin{longtable}{@{}l r r r r l@{}}
\caption{Proxy-level diagnostic of the current Level-B galactic
closure \eqref{eq:app_udg_closure} evaluated at
$R_{1/2}=\tfrac{4}{3}R_e$ using the isotropic Plummer Wolf
estimator.  Central values only; 95\% CL bands
in~\S\ref{app:udg:nuisance}.  AGC~114905 is a gas-rich rotator and
requires a disc kinematic solver rather than a spherical Jeans
proxy; it is listed for completeness of the UDG-class coverage.
DF2 remains observationally ambiguous.}
\label{tab:udg_stress_test}\\
\toprule
Object & $M_\star/10^8M_\odot$ & $R_e$/kpc &
$\sigma_{\rm obs}$/km\,s$^{-1}$ & $r=M_{\rm dyn}/M_{\rm bar}$ &
$\Xi_{\rm req}$\\
\midrule
\endfirsthead
\toprule
Object & $M_\star/10^8M_\odot$ & $R_e$/kpc &
$\sigma_{\rm obs}$/km\,s$^{-1}$ & $r=M_{\rm dyn}/M_{\rm bar}$ &
$\Xi_{\rm req}$\\
\midrule
\endhead
\bottomrule
\endfoot
NGC 1052--DF4        & 1.5  & 1.6  & $6.3^{+2.5}_{-1.6}$ & $\sim 1.1$ & $\sim 10^{-3}$ \\
FCC~224              & 1.74 & 1.89 & $7.8^{+6.7}_{-4.4}$ & $\sim 1.6$ & $\sim 2\times 10^{-2}$ \\
NGC 1052--DF2        & 1.3  & 2.2  & $[8.6,\,14.9]$      & $3$--$9$   & $0.1$--$1.4$ (ambiguous) \\
NGC~5846-UDG1        & 1.1  & 2.1  & $17\pm 2$           & $\sim 14$  & $\sim 1.6$ \\
Dragonfly 44         & 3.0  & 4.7  & $33\pm 3$           & $\sim 42$  & $\sim 8.9$ \\
AGC~114905 (rotator) & \multicolumn{5}{l}{Disc solver required; not included in this spherical diagnostic} \\
\end{longtable}

The structural observation is that within the six-object diagnostic
sample, the current closure at the baseline
$g^\dagger_{\rm eff}=g^\dagger_0$ returns $\Xi_{\rm req}\sim O(1)$
for normal RAR objects (NGC~5846-UDG1) and the upper tail
(Dragonfly~44), but $\Xi_{\rm req}\ll 1$ for NGC~1052--DF4 and
FCC~224 at their central observed dispersions.  NGC~1052--DF2
remains observationally ambiguous because its dispersion bracket
spans an order of magnitude across the
$\sigma\approx 8$--$15$~km\,s$^{-1}$ range discussed in the
literature.

\subsection{Matched-pair observation}
\label{app:udg:matched_pair}

The ECT-relevant tension is sharpest in the matched pair
NGC~1052--DF4 vs NGC~5846-UDG1: comparable $M_\star\sim 10^{8}M_\odot$,
comparable $R_e\sim 2$~kpc, comparable group-scale environment, but
central-value $\Xi_{\rm req}$ values differing by three orders of
magnitude.  Because the minimal density-only Level-B
ansatz~\eqref{eq:gdagger_env_main} depends only on the single
environmental scalar $\rho_{\rm env}$ and is monotonic in it, it
does not naturally accommodate both endpoints of such a matched
pair within a single simple environmental scalar law.  This is a
structural observation, not an impossibility theorem: the broader
ECT-compatible class of $\Xi$-functions noted
in~\S\ref{app:rar_derivation_status} remains open at OP3.

\subsection{Nuisance and uncertainty budget}
\label{app:udg:nuisance}

For DF4, propagating independent nuisances through
\eqref{eq:app_udg_xi_req} gives the diagnostic budget in
Table~\ref{tab:udg_nuisance}.  The multiplicative factors act on
$\Xi_{\rm req}$ for a fixed observational central value of
$\sigma_{\rm obs}$.

\begin{longtable}{@{}l l@{}}
\caption{Independent nuisance factors affecting the proxy-level
$\Xi_{\rm req}$ diagnostic for NGC~1052--DF4.  All factors are
multiplicative on $\Xi_{\rm req}$.  Combined 95\% CL:
$\log_{10}\Xi_{\rm req}\in[-4,\,-0.5]$.}
\label{tab:udg_nuisance}\\
\toprule
Source & Multiplicative factor on $\Xi_{\rm req}$\\
\midrule
\endfirsthead
\bottomrule
\endfoot
Distance $D\in[13,25]$~Mpc (covariant with $M_\star\propto D^2$, $R_e\propto D$) & $\sim 10\times$\\
Stellar mass-to-light ratio $M/L\in[1,4]$ (SPS-driven) & $\sim 10\times$\\
Small-$N$ Gaussian posterior at 95\% CL on $\sigma$ (7 GC tracers) & $\sim 100\times$\\
Anisotropy $\beta\in[-0.5,0.5]$ & $\sim 1.5\times$\\
Wolf-proxy $\to$ Jeans upgrade & $\sim 1.3\times$\\
\end{longtable}

The entire 95\% CL band for $\log_{10}\Xi_{\rm req}$ at DF4 lies
below the normal RAR band $\Xi\sim O(1)$.  The proxy-level
\emph{classification} of DF4 as a stress test of the current
closure is therefore robust to nuisance variations; only the
numerical value of $\Xi_{\rm req}$ is imprecise.  The present
bookkeeping is diagnostic-level: a posterior $P(\Xi\mid\text{data})$
with joint sampling of the full nuisance vector
$\theta=\{M/L,D,\beta,n,R_{\rm ap},f_{\rm int},M_{\rm host},d_{\rm 3D},Q_{\rm eq}\}$
is the appropriate next step.

\subsection{External field effect benchmark}
\label{app:udg:efe}

A benchmark parallel to the ECT diagnostic is the external field
effect (EFE) characteristic of MOND-class theories.  Estimating
$g_{\rm ext}\simeq GM_{\rm host}/d_{\rm proj}^2$ for the four
pressure-supported reference objects:

\begin{longtable}{@{}l c c c@{}}
\caption{External field amplitudes for the four spherical UDG
reference objects, for the MOND-EFE benchmark.  $a_0=1.2\times
10^{-10}\,\text{m s}^{-2}$ is the standard MOND acceleration
scale.  All four values lie far below $a_0$; EFE screening is
minor.}
\label{tab:udg_efe}\\
\toprule
Object & $M_{\rm host}/M_\odot$ & $d_{\rm proj}$/kpc & $g_{\rm ext}/a_0$\\
\midrule
\endfirsthead
\bottomrule
\endfoot
NGC 1052--DF4        & $\sim 10^{11}$   & $\sim 80$  & $\sim 0.02$\\
FCC~224              & $\sim 3\times 10^{10}$ & $\sim 70$  & $\sim 0.007$\\
NGC~5846-UDG1        & $\sim 5\times 10^{11}$ & $\sim 100$ & $\sim 0.06$\\
Dragonfly~44 (Coma)  & $\sim 10^{13}$   & $\sim 500$ & $\sim 0.05$\\
\end{longtable}

EFE screening is therefore minor for all four objects and, within
this benchmark, does not numerically distinguish DF4 from the
matched DM-rich control NGC~5846-UDG1.  This observation is
benchmark-level: the current ECT closure
\eqref{eq:app_udg_closure} does not carry an explicit
$g_{\rm ext}$ dependence, and the MOND-EFE comparison is strictly
outside the ECT formulation.

\subsection{Side diagnostics and sanity checks}
\label{app:udg:side_checks}

\paragraph{Ultra-faint dwarfs (side diagnostic, not validation).}
Applied to Milky-Way dwarf spheroidals at
$M_\star\sim 10^3$--$10^7 M_\odot$ (Draco, Sculptor, Fornax,
Segue~I, Willman~I), the same Wolf-based proxy returns
$\Xi_{\rm req}\in[\sim 2,\sim 10^3]$.  This is not a validation of
the proxy chain: the spread is driven by well-documented effects
orthogonal to the UDG stress test (binary-star inflation of
$\sigma_{\rm obs}$, tidal non-equilibrium, radial anisotropy, and
small-$N$ statistics) and merely confirms that the proxy does not
collapse outside the UDG target mass range.

\paragraph{Globular clusters (Newton sanity check).}
For 47~Tuc, $\omega$~Cen, M13, M15, the ratio
$g_N/g^\dagger_0\gtrsim 15$ places these systems safely in the
$g_N\gg g^\dagger_{\rm eff}$ branch of
\eqref{eq:app_udg_closure}, where the closure returns
$g\to g_N$ independently of $\Xi$.  This is the trivial Newton
recovery check at high density.

\paragraph{FCC~224 non-sphericity.}
FCC~224 is prolate with observed apparent ellipticity
$\varepsilon\simeq 0.36$ \cite{Buzzo2025FCC224}.  Spherical
Jeans should therefore be treated only as an order-unity proxy
for this object; no reliable numerical non-spherical correction
factor can be claimed within the present diagnostic.  The central
$\Xi_{\rm req}$ value for FCC~224 in
Table~\ref{tab:udg_stress_test} should accordingly be read as
order-of-magnitude.  This does not change its qualitative
classification as a stress-test object.

\paragraph{Selection rate.}
Approximately 3 confirmed DM-deficient UDGs among $\sim 500$
screened dwarfs with kinematic follow-up translates to
$\sim 0.6\%$ incidence \cite{Moreno2022DMfree}, consistent in
order of magnitude with the $\sim 1\%$ rate of collision-stripped
systems predicted by $\Lambda$CDM bullet-dwarf
simulations~\cite{Ogiya2018DF2tidal,vanDokkum2022bullet}.  Any
future ECT-compatible rescue mechanism will need to be consistent
with this observed rate.

\subsection{Selection and sample caveats}
\label{app:udg:sample}

The presently known DM-deficient UDG sample is small and
\emph{not} selection-neutral: several candidates were identified
through unusual globular-cluster-population markers
(top-heavy GCLF, excess luminous GCs for the host stellar mass).
The class properties inferred from this sample should therefore be
treated as provisional rather than as population-level conclusions.
Whether the class persists under selection-neutral future
surveys is itself part of the discriminant content of ND-13 in
Table~\ref{tab:future_discriminants}.

\subsection{Comparison with the $\Lambda$CDM bullet-dwarf interpretation}
\label{app:udg:lcdm}

In the $\Lambda$CDM + bullet-dwarf scenario
\cite{Ogiya2018DF2tidal,Moreno2022DMfree,vanDokkum2022bullet},
the DM-deficient UDGs arise as post-collisional remnants of
high-speed galaxy encounters in which the dark halos of the two
progenitors decouple from the stars and proceed approximately
ballistically, leaving two stellar remnants embedded in
essentially baryon-only potentials.  This scenario reproduces the
observed low dispersions through partial dark-halo stripping,
predicts a small-but-non-zero incidence rate compatible with the
$\sim 1\%$ screened-population observation, and places at least one
candidate (the compact dwarfs near NGC~1052 \cite{vanDokkum2022bullet})
in the same post-collision alignment.

The $\Lambda$CDM comparator therefore sits on fully derived
astrophysics, not on open elements of the underlying theory.  Any
future ECT-compatible rescue will need to reach comparable
explanatory power --- at minimum, reproduce the central dispersion
values of NGC~1052--DF4 and FCC~224 while remaining consistent
with the matched control NGC~5846-UDG1 and with the standard RAR
across the wider UDG population --- before it can be promoted to a
preferred interpretation within ECT.

\subsection{Possible theory-level directions for a future completion}
\label{app:udg:directions}

The full ECT field equations \eqref{eq:gen_EE} contain several
structural degrees of freedom not yet propagated through to the
practical galactic closure \eqref{eq:app_udg_closure}.  The
following are candidate directions for a future completion that
could in principle address the DM-deficient UDG class.  None is
presently derived for the galactic closure; each is listed here
only as a label for future work.  The list is neutral with respect
to which direction is more promising.
\begin{enumerate}
\item \textbf{Dynamical $G_{\rm eff}(\phi)$} in~\eqref{eq:gen_EE}.
  The generalised field equations contain a dynamical
  $G_{\rm eff}(\phi)$ entering through the $e^{\beta\phi}R$
  coefficient of the action~\eqref{eq:phi_first_action}.  A derivation
  mapping that degree of freedom onto the practical closure
  \eqref{eq:app_udg_closure} --- so that local disturbances of
  $\bar\phi$ appear in the observable $\sigma$ --- is not provided
  in the present preprint.  The route is structurally conceivable
  but not a derived result.
\item \textbf{History-dependent extensions of $\Xi$} within the
  broader ECT-compatible class described
  in~\S\ref{app:rar_derivation_status}.  The minimal
  ansatz~\eqref{eq:gdagger_env_main} depends only on the
  instantaneous environmental density.  An extension allowing
  $\Xi$ to depend on the local history of the condensate (e.g.\
  elapsed time since a relaxation-disrupting event) is within the
  open class at OP3, but no explicit form is constructed here.
\item \textbf{Orientation-sector coupling} through
  $\alpha_{\rm ECT}(n_\mu n_\lambda R^\lambda_\nu+\ldots)$ of
  \eqref{eq:gen_EE}.  An orientation-sensitive modification of
  the galactic closure is available in principle from the
  broader theory, but is not carried out here.
\item \textbf{Non-quasi-static $\phi$-dynamics} breaking the
  quasi-static assumption
  of~\eqref{eq:phi_eom_local}.  If the local condensate
  response time were comparable to the galactic dynamical time,
  the practical closure would acquire an additional
  time-dependent branch not captured by \eqref{eq:app_udg_closure};
  demonstrating this regime on the galactic scale requires a
  derivation that is not in the present preprint.
\end{enumerate}

The present appendix records the diagnostic; it does not promote
any of these routes to a solved intra-ECT rescue mechanism.

\subsection{Physical interpretation of the diagnostic}
\label{app:udg:interpretation}

Within the language of the current practical galactic closure, the
stress test has the following minimal physical reading.  The
closure \eqref{eq:app_udg_closure} admits two asymptotic branches:
Newtonian at $g_N\gg g^\dagger_{\rm eff}$ and deep-MOND at
$g_N\ll g^\dagger_{\rm eff}$.  In the DM-deficient UDGs the local
$g_N$ is small enough to enter the deep-MOND branch at the
baseline $g^\dagger_{\rm eff}=g^\dagger_0$, but the
\emph{observed} kinematics are Newton-like.  Within the scope of
the present closure, the only way to describe this within the
same single closure and the same asymptotic structure is to
interpret the closure as requiring a sharply suppressed local
$g^\dagger_{\rm eff}$ at these objects relative to the global
baseline $g^\dagger_0$.  This is the content of the diagnostic
$\Xi_{\rm req}\ll 1$ inferred above.  Whether such a suppression
arises from any of the directions
in~\S\ref{app:udg:directions} is an open question.

\subsection{Open questions tied to OP3}
\label{app:udg:open_questions}

The diagnostic of this appendix refines and sharpens the open
problem OP3 (exact environment law for $g^\dagger$) listed in the
main text.  Concretely:

\begin{itemize}
\item OP3a: Construct or exclude the density-only ansatz class.  Is
  there \emph{any} choice of $\Xi(\rho_{\rm env})$ with
  monotonically decreasing behaviour in the group-density regime
  that simultaneously reproduces matched DM-rich and DM-deficient
  UDGs at the observed levels?  The matched-pair observation
  in~\S\ref{app:udg:matched_pair} suggests not, but does not rule
  it out formally.
\item OP3b: Specify the additional variable.  If OP3a closes
  negatively, which of the directions
  in~\S\ref{app:udg:directions} supplies the additional
  environment/history variable that completes $\Xi$?
\item OP3c: Derive the galactic-scale effective action.  Perform
  the derivation that connects the dynamical
  $G_{\rm eff}(\phi)$ of \eqref{eq:gen_EE} to the practical
  galactic closure \eqref{eq:app_udg_closure}; only after this
  derivation can $G_{\rm eff}$-based effects on $\sigma_{\rm obs}$
  be claimed as ECT-internal rather than parametric.
\item OP3d: Predict the incidence rate.  Any completion that
  accommodates DF4 and FCC~224 must also be consistent with the
  observed $\sim 0.6\%$--$1\%$ incidence of DM-deficient UDGs in
  the screened dwarf population.
\end{itemize}

\subsection{What this appendix does and does not establish}
\label{app:udg:status}

\emph{Established here:}
\begin{itemize}
\item The practical galactic closure \eqref{eq:app_udg_closure} has
  exactly two asymptotic branches (Newton, deep-MOND); no third
  low-$g$ Newtonian branch exists within it.
\item Applied to normal RAR galaxies the diagnostic returns
  $\Xi_{\rm req}\sim O(1)$.
\item Applied to DF4 and (probably) FCC~224 the diagnostic returns
  $\Xi_{\rm req}\ll 1$ robustly under the nuisance budget of
  Table~\ref{tab:udg_nuisance}.
\item The minimal density-only Level-B $\Xi$-ansatz is
  structurally insufficient for the DF4 vs NGC~5846-UDG1 matched
  pair.
\end{itemize}

\emph{Not established here:}
\begin{itemize}
\item DF4 and FCC~224 are \emph{not} explained by the current
  closure.
\item No rescue mechanism is derived.  In particular,
  $G_{\rm eff}(\phi)$ is listed as a candidate direction on equal
  footing with the other items in~\S\ref{app:udg:directions}; it
  is not promoted to a solved intra-preprint mechanism.
\item The proposition
  ``$\bar\phi_{\rm local}\to 0^-\Rightarrow g^\dagger_{\rm eff}\to 0
  \Rightarrow$ Newton'' is \emph{not} a derived consequence of the
  attached preprint and must not be read into the above
  diagnostic.
\item The matched-pair comparison in
  \S\ref{app:udg:matched_pair} is a structural observation, not a
  formal impossibility theorem about all ECT-compatible
  $\Xi$-functions.
\end{itemize}

The integration target of this appendix is therefore the
discriminant entry ND-13 of Table~\ref{tab:future_discriminants},
not a claim of explanatory success.

% ============================================================
